\chapter{Relativistic one-electron layers: scalar and two-component X2C}
\label{ch:relativistic-layers}

\section{Purpose and scope}

The release carries two density-fitted X2C layers that differ in what
they transform and in what wavefunction they support.  The scalar
layer, \texttt{core.x2c\_casscf\_df}, applies the spin-free
Foldy--Wouthuysen one-electron Hamiltonian and runs a real
spin-adapted CASSCF on top of it.  The two-component layer,
\texttt{core.x2c\_df}, keeps the spinor structure of the X2C transform
and runs a complex generalized Hartree--Fock calculation in the
contracted spinor basis.

The two are not interchangeable.  The scalar layer has orbital and CI
response but no spin--orbit coupling; the two-component layer has
spin--orbit coupling at the one-electron level and no multireference
response.  Both use the non-relativistic Coulomb operator, density
fitted, with no picture-change correction on the two-electron part.
This is one-electron X2C (1eX2C); the separate \texttt{core.x2c\_sap}
module implements the atomic-potential approximation to the missing
two-electron picture-change term and is outside the validation record
here.

The scope is full-electron basis sets.  Both X2C helpers raise
\texttt{NotImplementedError} for an effective core potential, and the
layers inherit that limit.  The scalar layer is sized for small and
medium active spaces; the two-component layer is a Hartree--Fock layer.

\section{Theory and derivation}
\label{sec:x2c-theory}

\subsection{The four-component problem in block form}

The Dirac--Coulomb Hamiltonian in a molecular basis, after removing the
rest energy, has the block form
\begin{equation}
  H_{4c} =
  \begin{pmatrix}
    V & T\\
    T & \dfrac{W}{4c^{2}}-T
  \end{pmatrix},
  \qquad
  S_{4c} =
  \begin{pmatrix}
    S & 0\\
    0 & \dfrac{T}{2c^{2}}
  \end{pmatrix},
  \label{eq:x2c-4c}
\end{equation}
where, in the decontracted spinor basis used by the transform,
\begin{equation}
  S = \langle \chi_i|\chi_j\rangle,
  \quad
  T = \langle \chi_i|T_{\mathrm{kin}}|\chi_j\rangle,
  \quad
  V = \langle \chi_i|V_{\mathrm{nuc}}|\chi_j\rangle,
  \quad
  W = \langle \chi_i|\boldsymbol{\sigma}\cdot\mathbf{p}\,
      V_{\mathrm{nuc}}\,
      \boldsymbol{\sigma}\cdot\mathbf{p}|\chi_j\rangle.
  \label{eq:x2c-integrals}
\end{equation}
The generalized eigenproblem
\begin{equation}
  H_{4c}\,C = S_{4c}\,C\,E
  \label{eq:x2c-eigen}
\end{equation}
has positive- and negative-energy solutions.  Partition the
positive-energy eigenvectors into large and small components,
\begin{equation}
  C = \begin{pmatrix} C_L\\ C_S\end{pmatrix},
  \qquad
  X = C_S C_L^{-1}.
  \label{eq:x2c-x}
\end{equation}
The matrix $X$ is the small-to-large ratio of the positive-energy
branch.  The implementation obtains it from the eigenproblem rather
than solving the quadratic equation for $X$: it builds
Eq.~\eqref{eq:x2c-4c} with the spinor integrals, calls
\texttt{scipy.linalg.eigh}, selects the positive-energy block, and
forms $X$ by Eq.~\eqref{eq:x2c-x}.  This is the construction in
\texttt{core.x2c\_df.X2CTransform}.

\subsection{Block decoupling and the Foldy--Wouthuysen Hamiltonian}

Eliminating the small component with $X$ and retaining terms through
$\mathcal{O}(c^{-2})$ gives the effective two-component Hamiltonian
\begin{equation}
  S_{+} = S + \frac{1}{2c^{2}}X^{\dagger}TX,
  \label{eq:x2c-splus}
\end{equation}
\begin{equation}
  H_{+} = V + TX + (TX)^{\dagger} - X^{\dagger}TX
        + \frac{1}{4c^{2}}X^{\dagger}WX.
  \label{eq:x2c-hplus}
\end{equation}
The large-component metric is not the original overlap, so a
renormalization matrix $R$ is chosen such that
\begin{equation}
  R^{\dagger}S_{+}R = S,
  \qquad
  H_{2c} = R^{\dagger}H_{+}R.
  \label{eq:x2c-renorm}
\end{equation}
The implementation computes $R$ from the symmetric eigendecomposition
of $S$ and the transformed metric (the \texttt{\_metric\_renormalization}
helper).  The result is a complex Hermitian Hamiltonian in the
contracted spinor basis.  For a closed-shell determinant its
spin--orbit content is quenched at first order, but the operator still
differs from the spin-free one through the spin-dependent
renormalization and the spinor structure of $X$.

The scalar (spin-free) layer performs the same construction with
spin-averaged integrals: the transform is applied to the spatial
$S,T,V,W$ blocks, the resulting Hamiltonian is real and acts on spatial
orbitals, and the spin--orbit terms are omitted.  This is the
\texttt{SpinFreeX2CHelper} limit used by
\texttt{core.x2c\_casscf\_df}.

\subsection{Why the scalar and two-component energies differ}

The two shipped layers are not a controlled one-parameter comparison:
the scalar layer is CASSCF and the two-component layer is HF.  The
X2C-specific difference can be isolated only at the same wavefunction
level.  At that level there are three sources:

\begin{itemize}
\item the spin-free transform drops the spin-dependent part of the
  renormalization and of the picture-change matrix;
\item the two-component transform retains the spin--orbit operator,
  whose expectation value vanishes for a non-degenerate closed-shell
  determinant but whose second-order and response contributions do not;
\item the two-electron treatment differs when the scalar layer uses
  spatial DF factors and the two-component layer uses spinor RI-JK
  factors, even before correlation is added.
\end{itemize}
For a closed-shell singlet the first source usually dominates; for an
open-shell or degenerate system the second can dominate.  A comparison
that does not hold the wavefunction method and the two-electron
operator fixed measures all three together.

\subsection{The two-electron part and picture change}

Both layers use the non-relativistic Coulomb interaction,
\begin{equation}
  (pq|rs)=\iint
  \frac{\phi_p^{*}(\mathbf r)\phi_q(\mathbf r)\,
        \phi_r^{*}(\mathbf r')\phi_s(\mathbf r')}
       {|\mathbf r-\mathbf r'|}
  \,d\mathbf r\,d\mathbf r',
  \label{eq:x2c-coulomb}
\end{equation}
density fitted as in Chapter~\ref{ch:ri-df}.  The X2C transformation of
the one-electron operator induces a picture-change correction on the
two-electron operator; it is not applied here.  The missing term is the
dominant residual of 1eX2C for heavy elements, and the SAP-X2C module
is the separate approximation to it.

\subsection{The SOMF operator}

The Breit--Pauli spin--orbit operator for electron $i$ is
\begin{equation}
  H_{\mathrm{SO}}
  = \frac{\alpha^{2}}{2}\sum_i\left[
      \sum_A \frac{Z_A}{r_{iA}^{3}}
      \left(\mathbf r_{iA}\times\mathbf p_i\right)\cdot\mathbf s_i
      - \sum_{j\neq i}\frac{1}{r_{ij}^{3}}
      \left(\mathbf r_{ij}\times\mathbf p_i\right)
      \cdot\left(\mathbf s_i+2\mathbf s_j\right)
    \right],
  \label{eq:x2c-breit-pauli}
\end{equation}
with $\mathbf r_{iA}=\mathbf r_i-\mathbf R_A$,
$\mathbf r_{ij}=\mathbf r_i-\mathbf r_j$, and
$\alpha$ the fine-structure constant.  The first term is the nuclear
one-electron SOC; the second contains spin-same-orbit and
spin-other-orbit contributions.

The spin-orbit mean-field (SOMF) approximation replaces the
two-electron term by its contraction with the reference density $P$,
\begin{equation}
  h^{2e}_k[pq]
  = \sum_{rs}P_{rs}\left[
      (pq|g_k|rs)
      - \frac{3}{2}(pr|g_k|sq)
      - \frac{3}{2}(sq|g_k|pr)
    \right],
  \label{eq:x2c-somf}
\end{equation}
with the geometric kernel
\begin{equation}
  g_k(\mathbf r_i,\mathbf r_j)
  = \frac{\left(\mathbf r_{ij}\times\mathbf p_i\right)_k}
         {r_{ij}^{3}},
  \label{eq:x2c-somf-kernel}
\end{equation}
in the raw \texttt{int2e\_p1vxp1} convention.  The global sign of the
raw integral is the same as that of \texttt{int1e\_pnucxp}; it is a
convention and cancels in the ratio tests.  The coefficients $1$ and
$-3/2$ are the spin-summed spin-same-orbit and spin-other-orbit
combinations.  The factor $\alpha^{2}/2$ is applied once, in
Eq.~\eqref{eq:x2c-somf-total}: the \texttt{int2e\_p1vxp1} integral
carries the geometric kernel without it, and the raw one-electron
integral \texttt{int1e\_pnucxp} likewise carries no prefactor.  The
\texttt{ECPso} integral is already an effective operator in energy
units and receives no $\alpha^{2}/2$.  The one-electron nuclear term is
\begin{equation}
  h^{\mathrm{nuc}}_k[pq]
  = \sum_A Z_A
    \left\langle p\left|
      \frac{\left(\mathbf r_{iA}\times\mathbf p_i\right)_k}{r_{iA}^{3}}
    \right|q\right\rangle,
  \label{eq:x2c-nuclear-soc}
\end{equation}
where the sum runs over all centers for the all-electron operator; in
the ECP variant of Eq.~\eqref{eq:x2c-somf-ecp} it runs only over
centers without a spin--orbit ECP.  The SOMF operator is
\begin{equation}
  h^{\mathrm{SOMF}}_k
  = \frac{\alpha^{2}}{2}\left[h^{\mathrm{nuc}}_k+h^{2e}_k[P]\right].
  \label{eq:x2c-somf-total}
\end{equation}
The two-electron kernel is not symmetric under the pair exchange
$(pq)\leftrightarrow(rs)$, so the derivative of
$\mathrm{tr}(D\,h^{2e}[P])$ with respect to $P$ is not
$h^{2e}[D]$; it is the separate contraction in
\texttt{somf\_density\_response}.  A gradient of an SOC matrix element
must use that response, or the equivalent bilinear form.

\subsection{The ECP spin--orbit term}

With an effective core potential, the core electrons are absent from
the calculation and the spin--orbit part of the ECP replaces their
contribution.  On a center carrying a spin--orbit ECP the valence-level
operator is
\begin{equation}
  h^{\mathrm{SOMF\text{-}ECP}}_k
  = \frac{\alpha^{2}}{2}\left[h^{\mathrm{nuc}}_k+h^{2e}_k[P]\right]
    - V^{\mathrm{ECP}}_k,
  \label{eq:x2c-somf-ecp}
\end{equation}
where $h^{\mathrm{nuc}}_k$ is kept only on centers without a
spin--orbit ECP, $h^{2e}_k[P]$ is the valence $J-\tfrac{3}{2}K$ mean
field of the spin-summed density $P$ with
$\mathrm{tr}(PS)=N_{\mathrm{elec}}$, and $V^{\mathrm{ECP}}_k$ is the
three-component \texttt{ECPso} integral.  On a center carrying a
spin--orbit ECP the ECP term replaces the nuclear term: the ECP is
fitted to reproduce the all-electron spin--orbit splitting of the
valence configuration and already contains the nuclear one-electron
term and the core--valence screening for that center.  The
Breit--Pauli terms carry $\alpha^{2}/2$; \texttt{ECPso} is already an
effective operator in energy units and is subtracted.

There is no unique first-principles decomposition of an ECP into one-
and two-electron spin--orbit parts.  The ECP is fitted to atomic
reference states, so part of the valence--valence mean field is already
folded into its parameters.  Adding the explicit valence
$J-\tfrac{3}{2}K$ term is therefore an approximation with a
double-counting error that cannot be removed without the ECP generation
data.  The pieces are exposed independently
(\texttt{nuclear=}, \texttt{ecp\_term=}) so variants can be compared.

\section{Input and operation}
\label{sec:x2c-input}

\subsection{Common conventions}

\begin{itemize}
\item \texttt{mol} is a built \texttt{pyscf.gto.Mole}.  The scalar layer
  expects the spatial orbital basis; the two-component layer builds its
  own decontracted spinor molecule internally through
  \texttt{SpinorX2CHelper} and never asks the caller for a spinor
  basis.
\item Energies are in Eh.  Nuclear gradients are in Eh/Bohr.  Atom
  coordinates may be supplied in Angstrom or Bohr through the usual
  pyscf \texttt{unit} argument; the gradient is always Eh/Bohr.
\item Only full-electron basis sets are accepted.  Both X2C helpers
  raise \texttt{NotImplementedError} for a molecule that carries an
  effective core potential.
\item The auxiliary basis defaults to \texttt{def2-universal-jkfit},
  resolved by \texttt{core.ri.make\_auxmol}, which falls back to the
  orbital basis when a fitting set does not cover every element.  The
  label actually used is \texttt{result.auxbasis} for the scalar
  result and \texttt{result.factors.auxbasis} for the two-component
  result; it should be quoted with any number.
\item The scalar layer's \texttt{ncore}, \texttt{nact}, and
  \texttt{nele} are spatial-orbital and electron counts, not spinor
  counts.  The two-component layer has no active-space interface.
\end{itemize}

A minimal skeleton for both layers is
\begin{lstlisting}[language=Python]
import numpy as np
from pyscf import gto
from diffmrc.core import x2c_casscf_df as xc
from diffmrc.core import x2c_df

mol = gto.M(atom='O 0 0 0; H 0 -0.757 0.587; H 0 0.757 0.587',
            basis='6-31g', unit='Angstrom', verbose=0)

# Scalar (spin-free) CASSCF, CAS(2,2), four core spatial orbitals.
res = xc.df_x2c_casscf(mol, ncore=4, nact=2, nele=2,
                       auxbasis='def2-universal-jkfit')
grad, parts = xc.df_x2c_casscf_gradient(mol, res)

# Two-component Hartree-Fock on the same molecule.
h = x2c_df.df_x2c_hcore(mol)
hf = x2c_df.df_x2c_hf(mol, auxbasis='def2-universal-jkfit')
g_hf, parts_hf = x2c_df.df_x2c_hf_gradient(mol, result=hf)
\end{lstlisting}

\subsection{Scalar DF-X2C CASSCF}

The energy is
\begin{equation}
  E = \sum_{pq}\gamma_{pq}h_{pq}
    + \tfrac{1}{2}\sum_{pqrs}\Gamma_{pqrs}(pq|rs)+E_{\mathrm{nuc}},
  \qquad
  h=\texttt{SpinFreeX2CHelper}(mol).\texttt{get\_hcore}(mol),
  \label{eq:x2c-scalar-energy}
\end{equation}
with the density-fitted Coulomb operator of
Eq.~\eqref{eq:x2c-coulomb}.  The wavefunction is a real spin-adapted
determinant expansion over the active spin orbitals.

\begin{center}
\small
\begin{tabularx}{\textwidth}{>{\raggedright\arraybackslash}p{0.42\textwidth}X}
\toprule
Entry point & Returns and purpose \\
\midrule
\texttt{sfx2c\_hcore(mol)} &
  Spin-free X2C one-electron Hamiltonian in the spatial AO basis. \\
\texttt{sfx2c\_hcore\_derivatives(mol)} &
  Its nuclear derivatives, shape \texttt{(natm, 3, nao, nao)}. \\
\texttt{spatial\_df\_factors(mol, auxmol=None, auxbasis=None,
  rcond=1e-12)} &
  \texttt{SpatialDFFactors} with \texttt{b}, \texttt{s},
  \texttt{jhalf}, the metric eigendecomposition, the auxiliary label,
  and the factor derivatives. \\
\texttt{df\_x2c\_rhf(mol, factors=None, auxmol=None, auxbasis=None,
  conv\_tol=1e-12, max\_cycle=200, dm0=None)} &
  Closed-shell reference SCF on the same Hamiltonian. \\
\texttt{df\_x2c\_casscf(mol, ncore, nact, nele, auxmol=None,
  auxbasis=None, nroots=1, weights=None, target=0, max\_cycle=300,
  g\_conv=1e-6, rcond=1e-12, theta0=None, rhf=None, verbose=False)} &
  \texttt{DFX2CCASSCFResult} with orbitals, CI vectors, state
  energies, state-averaged energy, and orbital-gradient norm. \\
\texttt{df\_x2c\_casscf\_gradient(mol, result=None, state=None,
  auxmol=None, auxbasis=None, rcond=1e-12)} &
  \texttt{(grad, parts)} for one CI root, shape \texttt{(natm, 3)} in
  Eh/Bohr; parts are one-electron, two-electron, Pulay, nuclear. \\
\texttt{df\_x2c\_energies\_at(mol, result, theta=None)} &
  State energies at a new geometry with the rotation parameters fixed;
  the reference orbitals follow the geometry through a displaced SCF
  aligned to \texttt{result.c\_ref}. \\
\texttt{df\_x2c\_sa\_state\_gradient(mol, result=None, state=None,
  fd\_step=1e-3, theta\_step=1e-5, rcond=1e-8, auxmol=None,
  auxbasis=None, return\_info=False)} &
  State-specific gradient of a state-averaged calculation; requires
  \texttt{result.nroots >= 2}. \\
\bottomrule
\end{tabularx}
\end{center}

\texttt{nroots=1} is the single-state calculation.  With
\texttt{nroots>1}, \texttt{weights} default to equal weights and
\texttt{target} selects the state reported in \texttt{result.e\_tot}.
The state-averaged energy is \texttt{result.e\_sa}; the individual
energies are \texttt{result.state\_energies}.  The \texttt{converged}
flag is \texttt{gnorm < 1e-5}; the energy and the nuclear gradient are
insensitive to the residual along flat orbital directions through the
envelope theorem.

The single-state gradient is
\begin{equation}
  \frac{dE}{dR}
  = \sum_{pq}\gamma_{pq}\frac{dh_{pq}}{dR}
  + \sum_P\sum_{pq}\frac{dw_{pqP}}{dR}X_{pqP}
  - \sum_{pq}F_{pq}\frac{dS_{pq}}{dR}
  + \frac{dE_{\mathrm{nuc}}}{dR},
  \qquad
  X_{pqP}=\sum_{rs}w_{rsP}\Gamma_{pqrs}.
  \label{eq:x2c-scalar-gradient}
\end{equation}
The one-electron derivative comes from pyscf's spin-free X2C hcore
gradient generator; the three-center and metric derivatives reuse
Chapter~\ref{ch:ri-df}.  The gradient entry point requires the
converged \texttt{result}.

A state-averaged calculation is different: the orbitals minimize
$E_{\mathrm{SA}}$, so the target-state orbital gradient $g_t$ does not
vanish, and
\begin{equation}
  \frac{dE_t}{dR}
  = \left.\frac{dE_t}{dR}\right|_{\theta}+z^{T}M,
  \qquad
  Hz=-g_t,
  \qquad
  H=\frac{\partial^{2}E_{\mathrm{SA}}}{\partial\theta^{2}},
  \label{eq:x2c-scalar-sa}
\end{equation}
with $M=\partial g_{\mathrm{SA}}/\partial R$.  Both $H$ and $M$ are
finite-difference constructions around the analytic state-averaged
orbital gradient: $H$ over $\theta$ (the CI vectors are re-diagonalized
at every step, so the CI response is included) and $M$ over displaced
geometries, rebuilding the integrals and the geometry-dependent
reference orbitals.  The direct term is a central difference of the
target-state energy along the envelope path; the closed-form
single-state formula does not apply because the target state is not
stationary at the SA orbitals.  The Hessian is pseudo-inverted with a
relative cutoff; the retained directions are reported in
\texttt{info}.

\subsection{Two-component DF-X2C Hartree--Fock}

The two-component layer solves
\begin{equation}
  E=\mathrm{tr}(DH_{\mathrm{X2C}})
    +\tfrac{1}{2}\mathrm{tr}(DJ)
    -\tfrac{1}{2}\mathrm{tr}(DK)+E_{\mathrm{nuc}},
  \label{eq:x2c-two-component-energy}
\end{equation}
with the RI-JK factorization
\begin{equation}
  (ij|kl)\approx\sum_P w_{ijP}w_{klP},
  \qquad
  w=BS^{-1/2},
  \qquad
  B_{ijP}=(ij|P),
  \label{eq:x2c-two-component-df}
\end{equation}
built directly with \texttt{int3c2e\_spinor}.  The one-electron
operator is Eq.~\eqref{eq:x2c-renorm} from the spinor integrals.  The
values agree with pyscf's \texttt{SpinorX2CHelper} to
$4.5\times10^{-11}$ on H$_2$O/6-31G (test bound $10^{-9}$ absolute and
$10^{-10}$ relative).

The spinor ERI convention is the complex-GHF one,
\begin{equation}
  J_{ij}=\sum_{kl}E_{ijkl}D_{lk},
  \qquad
  K_{ij}=\sum_{kl}E_{ikjl}D_{kl},
  \label{eq:x2c-spinor-jk}
\end{equation}
which for $w=BS^{-1/2}$ becomes
\begin{equation}
  \rho_P=\sum_{ij}w_{ijP}D_{ji},
  \qquad
  J=\sum_Pw_P\rho_P,
  \qquad
  K=\sum_Pw_PDw_P.
  \label{eq:x2c-spinor-jk-factor}
\end{equation}
The three-center tensor is Hermitian in the AO pair,
$B_{ijP}=\overline{B_{jiP}}$, so the derivative with respect to the
second AO center carries a complex conjugate.  The conjugate was
verified by finite differences of $B$ at $2.5\times10^{-9}$ relative
on H$_2$O/6-31G.

\begin{center}
\small
\begin{tabularx}{\textwidth}{>{\raggedright\arraybackslash}p{0.44\textwidth}X}
\toprule
Entry point & Returns and purpose \\
\midrule
\texttt{df\_x2c\_hcore(mol, C=None, c=LIGHT\_SPEED,
  return\_deriv=False, step=1e-4)} &
  Complex Hermitian two-component Hamiltonian, optionally in a
  supplied basis and optionally with \texttt{dH/dR}, shape
  \texttt{(natm, 3, n, n)}. \\
\texttt{X2CTransform(mol, c=LIGHT\_SPEED)} &
  Transform at one geometry: decontracted spinor molecule, spinor
  contraction matrix, and the four input integrals. \\
\texttt{spinor\_df\_factors(mol, auxmol=None, auxbasis=None,
  rcond=1e-12)} &
  \texttt{SpinorDFFactors} with \texttt{b}, \texttt{s}, \texttt{jhalf},
  metric eigendecomposition, and factor derivatives. \\
\texttt{df\_x2c\_hf(mol, auxmol=None, auxbasis=None, dm0=None,
  conv\_tol=1e-11, max\_cycle=200, rcond=1e-12, verbose=False)} &
  \texttt{DFX2CHFResult} with spinor density, orbitals, orbital
  energies, occupations, Fock matrix, energy parts, and gradient norm. \\
\texttt{df\_x2c\_hf\_gradient(mol, result=None, auxmol=None,
  auxbasis=None, dm0=None, transform\_step=1e-4, rcond=1e-12,
  verbose=False)} &
  \texttt{(grad, parts)} for the converged DF-X2C HF energy. \\
\texttt{energy\_weighted\_density(result)} &
  $W=C_{\mathrm{occ}}\mathrm{diag}(\epsilon_i n_i)
   C_{\mathrm{occ}}^{\dagger}$. \\
\texttt{spinor\_overlap\_derivatives(mol)} &
  $dS/dR$ of the contracted spinor overlap. \\
\bottomrule
\end{tabularx}
\end{center}

At a converged SCF the density response collapses to the Pulay term:
\begin{equation}
  \frac{dE}{dR}
  = \sum_{ij}D_{ij}\frac{dH_{ji}}{dR}
  + \left.\frac{dE_{2e}}{dR}\right|_{D}
  - \sum_{ij}W_{ij}\frac{dS_{ji}}{dR}
  + \frac{dE_{\mathrm{nuc}}}{dR}.
  \label{eq:x2c-two-component-gradient}
\end{equation}
With $G_P=\rho_PD^{T}-(Dw_PD)^{T}$,
\begin{equation}
  dE_{2e}=\sum_{ijP}G_{ijP}\,dw_{ijP},
  \qquad
  dw=dBS^{-1/2}+B\,dS^{-1/2},
  \label{eq:x2c-two-component-dE}
\end{equation}
and the metric inverse-square-root derivative uses the
Daleckii--Krein form of Chapter~\ref{ch:ri-df}.  The X2C transform
derivative combines analytic nuclear derivatives of the four input
integrals, the implicit eigenbasis derivative of the picture-change
matrix (a complement-space sum safe at the exact Kramers degeneracy),
and a transform-layer central difference at $10^{-4}$ Bohr for the
polynomial and matrix-square-root chain; the development probes
measure its error at about $10^{-7}$ relative on H$_2$O.

\section{Worked example}
\label{sec:x2c-example}

\subsection{OH and OF: one-electron versus SOMF spin--orbit splitting}

The SOMF validation on OH and OF uses two state-averaged roots, the
same geometry and active space for the stack and ORCA 6.1.1, and the
SI convention of this stack.  The recorded numbers are
\begin{center}
\small
\begin{tabular}{lrrrr}
\toprule
System & one-electron (ours) & SOMF (ours) & ORCA one-electron & ORCA SOMF \\
\midrule
OH & $206.52$ & $129.23$ & $205.88$ & $128.83$ \\
OF & $273.06$ & $176.77$ & $272.06$ & $176.12$ \\
\bottomrule
\end{tabular}
\end{center}
in cm$^{-1}$.  The screening ratio SOMF/one-electron removes the
overall one-electron convention offset between the codes.  The
recorded ratios are $0.625738$ (OH) and $0.647358$ (OF), against ORCA
$0.625753$ and $0.647357$; the scale-corrected SOMF values are
$128.827$ and $176.120$~cm$^{-1}$, with deviations
$-0.0030$~cm$^{-1}$ ($-2.35\times10^{-5}$ relative) and
$+0.00024$~cm$^{-1}$ ($1.34\times10^{-6}$ relative).  The source file
is \texttt{experiments/somf/results\_somf.json}; the driver is
\texttt{experiments/somf/validate\_oh\_of.py}.

The one-electron column is not the physical SOC; it omits the
two-electron screening.  The SOMF column is the mean-field
approximation to that screening.  The near-perfect ratio agreement is
the evidence that the two-electron contraction has the right shape,
not that the absolute operator is exact.

\subsection{GeF with an ECP: the ECP spin--orbit operator}

The ECP validation on GeF uses \texttt{lanl2dz} on Ge and F, the
spin--orbit ECP \texttt{ecpds28mdfso} on Ge, and a two-state
state-averaged calculation.  The recorded valence-level SOMF splitting
is $28.37$~cm$^{-1}$, against $43.72$~cm$^{-1}$ for the nuclear-only
operator on the same density.  The difference is the valence mean
field and is the quantity the ECP decomposition does not fix uniquely.
On YbF the \texttt{ECPso} norm on Yb is $17.04$ per component, while
the raw nuclear operator norm on the same center is of order
$10^{4}$; adding both would be dominated by an operator that the ECP
was built to replace.  The gradient prototype on GeF follows a finite
difference of the SOC matrix element to $2.4\times10^{-5}$ relative,
with a density-response term of $4.4\times10^{-7}$ in that direction.
The source files are
\texttt{experiments/somf\_ecp/results\_compare\_orca\_gef.json} and
\texttt{experiments/somf\_ecp/results\_somf\_grad\_f.json}.

\subsection{H\texorpdfstring{$_2$}{2}O/6-31G: two-component X2C
acceptance bounds}

The two-component DF-X2C HF layer on H$_2$O/6-31G with
\texttt{def2-universal-jkfit} is checked against pyscf's
\texttt{SpinorX2CHelper} and exact-ERI X2C UHF.  The recorded
acceptance bounds are
\begin{center}
\small
\begin{tabular}{lr}
\toprule
Check & Bound \\
\midrule
hcore against \texttt{SpinorX2CHelper} & max $<10^{-9}$, relative $<10^{-10}$ \\
DF SCF against exact X2C UHF & $<10^{-4}$~Eh (the fit) \\
gradient against 3-point differences & relative $<2\times10^{-5}$ \\
gradient against 5-point differences & relative $<2\times10^{-6}$ \\
translation and rotation invariance & relative $<10^{-6}$ \\
factor derivatives against geometry FD & relative $<10^{-6}$ \\
\bottomrule
\end{tabular}
\end{center}
These are the release's recorded criteria; the raw gradient components
are not stored in the release tree.  A requested XeH$^{+}$
three-direction gradient record is also not part of this tree; the
XeH$^{+}$ file that is present
(\texttt{tests/core/test\_x2c\_sap\_heavy.py}) checks only that the
SAP-X2C energy moves toward the four-component DHF reference and is
gated behind \texttt{DIFFMRC\_X2C\_SAP\_HEAVY=1}.  The source files for
the table are \texttt{tests/core/test\_x2c\_df.py} and
\texttt{experiments/df\_x2c\_gradient\_validation.py}.

\section{Pitfalls}
\label{sec:x2c-limits}

\subsection{Scalar layer}

{\sloppy
\begin{itemize}
\item Effective core potentials are not supported.  The spin-free X2C
  helper raises \texttt{NotImplementedError} for an ECP molecule, and
  the layer inherits that limit.
\item The CI, RDM, and Fock contractions are dense in the spin-orbital
  space; the 2-RDM is an $n_{\mathrm{so}}^4$ object.  The prototype is
  sized for the validation systems.  Production f-block active spaces
  need a factorized 2-RDM route.
\item The complex two-component CI, complex orbital rotations, and the
  Kramers-restricted response are designed but not implemented.
\item The state-averaged response is a finite-difference construction
  of $H$ and $M$ around the analytic orbital gradient.  Its cost grows
  with the number of rotation parameters and displaced geometries.
\item The displaced-geometry reference has a second orbital minimum on
  H$_2$O in the same $\theta$ space.  \texttt{df\_x2c\_energies\_at}
  follows the envelope path; re-optimizing at the displaced geometry is
  a different, branch-sensitive path and is not the reference for the
  analytic gradient.
\end{itemize}
\par}

\subsection{Two-component layer}

\begin{itemize}
\item The layer is Hartree--Fock.  A state-specific or state-averaged
  two-component CASSCF is not implemented.
\item Full-electron basis sets only; the ECP limit is the same as in
  the scalar layer.
\item The physical X2C multireference path remains open.  The pyscf
  \texttt{int2e\_spinor} array does not obey the chemist pair symmetry
  (measured violation $0.62$ on H$_2$O/6-31G), and the CI Hamiltonian
  built from it with the determinant-energy ordering does not commute
  with many-body time reversal on OH/6-31G.  Kramers-restricted
  response is therefore not available on this stack.
\item The X2C transform derivative carries a transform-layer central
  difference at $10^{-4}$ Bohr, with a measured error about $10^{-7}$
  relative on H$_2$O.  It is not an analytic derivative of the
  transform chain.
\item A scalar-versus-two-component energy difference is not a pure
  X2C effect when the two layers use different wavefunction methods.
  Hold the method and the two-electron operator fixed before assigning
  the difference to relativity.
\end{itemize}

\subsection{SOMF and ECP}

\begin{itemize}
\item The SOMF operator is density dependent.  A gradient of an SOC
  matrix element must include the density response; the
  Hellmann--Feynman expression alone is incomplete.
\item The two-electron kernel is not symmetric under
  $(pq)\leftrightarrow(rs)$; use \texttt{somf\_density\_response},
  not the value contraction, for a derivative.
\item The ECP decomposition is not unique.  The split between the ECP
  spin--orbit operator and the explicit valence mean field is a fitting
  freedom; report the variant bracket.
\item Eligibility is a property of the ECP data, not of the atomic
  number.  A scalar-only ECP does not supply \texttt{ECPso}.
\end{itemize}

\section{Validation and references}
\label{sec:x2c-validation}

{\raggedright
The files behind this chapter are
\texttt{tests/core/test\_x2c\_casscf\_df.py},
\texttt{tests/core/test\_x2c\_df.py},
\texttt{tests/core/test\_somf.py},
\texttt{tests/core/test\_somf\_ecp.py},
\texttt{experiments/df\_x2c\_casscf\_validation.py},
\texttt{experiments/df\_x2c\_gradient\_validation.py},
\texttt{experiments/somf/validate\_oh\_of.py},
\texttt{experiments/somf\_ecp/validate\_ecp\_somf.py}, and
\texttt{experiments/somf\_ecp/validate\_somf\_grad.py}.  The measured
values are the OH/OF splittings, the GeF ECP splitting and gradient,
the YbF ECP norms, and the acceptance bounds of
Section~\ref{sec:x2c-example}.\par}

The X2C construction follows W.~Kutzelnigg and W.~Liu,
\emph{J. Chem. Phys.} \textbf{123}, 241102 (2005), and the review in
M.~Reiher and A.~Wolf, \emph{Relativistic Quantum Chemistry}
(Wiley, 2009).  The SOMF operator follows the ORCA 6.1 manual
(Section 5.26), F.~Hess, C.~M.~Marian, U.~Wahlgren and O.~Gropen,
\emph{Chem. Phys. Lett.} \textbf{251}, 365 (1996), and F.~Neese,
\emph{J. Chem. Phys.} \textbf{122}, 034107 (2005).  The
effective-core-potential construction and its spin--orbit integrals
are documented in the pyscf \texttt{gto} and \texttt{ECPso}
documentation and in the Stuttgart--K\"oln ECP papers cited by the
\texttt{stuttgart\_rsc} basis data.
